\chapter{Tag 1 — Prüfziffern und Fehlererkennung}

Stub-Kapitel für M1. Inhalte folgen in M4.

Dieses Kapitel testet \emph{Szenario 1} aus dem Layout-Briefing:
eine Aufgabe pro Aufgabentyp auf einer Seite, alle drei Symbole
sichtbar.

\begin{bleistiftuebung}{1}{Wiederholung als Code}
Welche der folgenden empfangenen Dreierblöcke werden korrekt zu
\(0\) oder \(1\) zurückübersetzt?
\texttt{111}, \texttt{010}, \texttt{001}, \texttt{110},
\texttt{000}, \texttt{101}.

Wie viele Bitfehler pro Dreierblock kann das Verfahren korrigieren?
\end{bleistiftuebung}

Ein bisschen Fließtext zwischen den Aufgaben, damit die Marginalien
sichtbar Abstand bekommen und das Zusammenspiel der Boxen geprüft
werden kann.

\begin{pythoneinheit}{1}{Paritätsprüfer bauen}
Lege in Thonny eine neue Datei an und tippe diesen Code ab — das
testet \emph{Szenario 3} (Code-Block innerhalb eines Aufgaben-Blocks):

\pythoncode{tag1/paritaet.py}

Speichere die Datei als \texttt{paritaet.py} und führe sie aus.
\end{pythoneinheit}

\begin{werkzeugcheck}{Thonny einrichten}
Lade Thonny von \texttt{thonny.org} herunter und installiere es. Du
findest dort ein Installationspaket für dein Betriebssystem.
\end{werkzeugcheck}
